\documentclass[9pt]{exam} 
\addpoints
\pointsinmargin

\usepackage[fleqn]{amsmath}
\usepackage{amssymb} %standard classy maths symbols
\usepackage[a4paper, margin=1in]{geometry} %makes it easy to specify where to put stuff on the page

\usepackage{tcolorbox} % basic but good package for boxes around text
\tcbset{colback=black!5!white}

\usepackage{tikz} % interval diagrams, sets, etc.

\begin{document}

\textbf{Name:} \dotfill

\

\begin{questions}

\question[2] From the definition, prove that $f_1(x) = a \in \mathbb{R}$ is continuous $\forall x \in \mathbb{R}$.
\fillwithdottedlines{20mm}

\question[2] From the definition, prove that $f_2(x) = x \in \mathbb{R}$ is continuous $\forall x \in \mathbb{R}$.
\fillwithdottedlines{20mm}


\question[2] Using theorems proved in class, explain why $f_2 - f_1$ is continuous on $\mathbb{R}$.
\fillwithdottedlines{20mm}

\question[2] Let $f$ be continuous on the closed interval $[a, b]$. Using theorems proved in class, explain why $g(x) = k \cdot (x - a) + f(a) - f(x)$ is also continuous on $[a, b]$ for $k \in \mathbb{R}$ constant.
\fillwithdottedlines{20mm}

\question[3] The function $g$ is important in the proof of which theorem? \newline
Find the statement of this theorem in the CRM and copy the diagram here.


\vfill

\question[0] \textbf{Bonus.} Using the rules, calculate the derivative of $h(x) = x \sin \big(\frac{1}{x}\big)$. \newline
Where is $h$ continuous? Where is $h$ differentiable? \hfill \emph{Work on the back.}

\end{questions}


\begin{tcolorbox}

\textbf{Name of marker:}

\begin{center}
\gradetable[h][questions]
\end{center}

\end{tcolorbox}


\end{document}